\documentclass[11pt]{article}

\usepackage[T1]{fontenc}
\usepackage{lmodern}
\usepackage{amsmath,amssymb,amsthm}
\usepackage[margin=1in]{geometry}
\usepackage{booktabs}
\usepackage{microtype}
\usepackage[hidelinks]{hyperref}

\newtheorem{theorem}{Theorem}
\newtheorem{proposition}[theorem]{Proposition}
\newtheorem{lemma}[theorem]{Lemma}
\newtheorem{corollary}[theorem]{Corollary}
\theoremstyle{remark}
\newtheorem{remark}[theorem]{Remark}

\newcommand{\rad}{\operatorname{rad}}
\newcommand{\cosoc}{\operatorname{cosoc}}
\newcommand{\ord}{\operatorname{ord}}
\newcommand{\QQ}{\mathbb{Q}}
\newcommand{\OK}{\mathcal{O}_K}

\title{Radicals in iterated quadratic \(abc\)-transfers}
\author{Magnus Gille\\Independent researcher}
\date{July 2026}

\begin{document}
\maketitle

\begin{abstract}
Let \((a_0,b_0,c_0)\) be a primitive positive triple with
\(a_0+b_0=c_0\), where \(a_0,b_0\) have opposite parity, and iterate the
classical quadratic transfer
\[
 (a,b,a+b)\longmapsto\bigl(4ab,(a-b)^2,(a+b)^2\bigr).
\]
Writing \(d_j=a_j-b_j\), \(R_n=\rad(a_nb_nc_n)\), and
\[
 W_n=\prod_{j<n}\frac{|d_j|}{\rad(|d_j|)},
\]
we prove an exact radical identity.  If
\(\cos\theta=(b_0-a_0)/c_0\), then
\[
 \frac{R_n}{c_n}
 =\frac{R_0}{c_0}\,
   \frac{|\sin(2^n\theta)|}
        {2^n\sin\theta\,W_n}.
\]
One explicit Baker--Wüstholz estimate gives
\[
 \log(c_n/R_n)=\log W_n+\Theta_{a_0,b_0}(n),
 \qquad \log c_n=2^n\log c_0.
\]
Thus every such orbit eventually consists of \(abc\)-hits, while a fixed
asymptotic quality gap is equivalent to positive-power growth of \(W_n\).
The orbit is also a dyadic subfamily of known negative-discriminant
Lucas \(abc\)-triples.  Our orbit-adapted argument recovers along dyadic
indices an archimedean limit already implied by general all-index Lucas
estimates.  For the seed \((1,8,9)\),
\[
 |d_j|=\frac12\bigl|V_{2^{j+1}}(2,9)\bigr|.
\]
Standard Lucas
valuation theory turns \(W_n\) into an aggregate of Lucas--Wieferich
contributions.  A reproducible search found no square prime divisor for
\(p\leq10^7\) and \(0\leq j\leq50\); the computation is not used in the
proofs.
\end{abstract}

\noindent\textit{2020 Mathematics Subject Classification.}
Primary 11B39; Secondary 11D75, 11J86, 11Y55.

\smallskip
\noindent\textit{Keywords.}
\(abc\) conjecture, radical, Lucas sequence, Chebyshev polynomial,
linear forms in logarithms, Wieferich prime.

\section{Introduction}

For a positive integer \(m\), write
\(\rad(m)=\prod_{p\mid m}p\).  A primitive triple of positive integers
\((a,b,c)\) with \(a+b=c\) is called an \(abc\)-hit if
\(\rad(abc)<c\).  Its quality is
\[
 q(a,b,c)=\frac{\log c}{\log\rad(abc)}.
\]
The \(abc\) conjecture predicts that, for every \(\varepsilon>0\), only
finitely many such triples have \(q(a,b,c)>1+\varepsilon\); see
Granville and Tucker~\cite{GranvilleTucker2002} for an introduction.

The identity
\[
 4ab+(a-b)^2=(a+b)^2                                      \tag{1.1}
\]
is a classical polynomial transfer for \(abc\)-triples.  It appears in
Oesterl\'e's discussion of the \(abc\)--Szpiro connection
\cite[p.~170]{Oesterle1988} and in the catalogue of Martin and
Miao~\cite[\S2.4]{MartinMiao2016}; see also
van der Horst~\cite[\S2.3]{vanderHorst2010}.  Iterating this exact
transfer to obtain infinitely many hits, together with one-step quality
control, was described by Holszty\'nski~\cite{Holsztynski2017} and
Guninski~\cite{Guninski2020}.  Thus neither the transfer nor its
iteration is new.  For odd \(c\), Ohana, Spicer, and
Stein~\cite[Prop.~1]{OhanaSpicerStein2013} also record the exact
one-step formula
\[
 \rad\bigl((b-a)^2(4ab)c^2\bigr)
 =\rad(|b-a|)\rad(abc).
\]
Martin and Miao~\cite[\S3.1]{MartinMiao2016} iterate a different
transfer and obtain one-sided radical and quality bounds.
Classical families already include
\((1,9^k-1,9^k)\)~\cite{AlvarezEtAl2023}, so neither the growth law
\(c_{n+1}=c_n^2\) nor the mere production of infinitely many hits is a
novelty claim here.

Our purpose is quantitative.  We first prove an exact disjoint
prime-support factorization along the whole orbit.  It separates the
radical loss into an archimedean product of Chebyshev values and a
nonarchimedean product \(W_n\) containing every repeated prime factor.
The first product telescopes to one sine value.  A single linear-form
estimate then bounds its logarithm above and below on the linear scale
\(n\), exponentially smaller than \(\log c_n\).  This gives an exact
radical identity, effective two-sided control, and a two-way equivalence
between the asymptotic \(abc\)-quality of the orbit and aggregate
repeated-prime growth.

These orbits also sit inside a Lucas-sequence construction already
studied by Bolvardizadeh~\cite[Ch.~2]{Bolvardizadeh2023}.  Her treatment
of quality assumes positive discriminant, and her list of future
directions isolates a missing archimedean limit when the discriminant is
negative~\cite[pp.~63--64]{Bolvardizadeh2023}.  We identify our orbit as
a constrained dyadic subfamily of that negative-discriminant setting and
give a direct, seed-specific estimate along dyadic indices.  The later
general theorem of Hajdu and Tijdeman~\cite[Theorem~2.1 and
Remark~1]{HajduTijdeman2025} gives non-real Lucas growth estimates for
all indices, and therefore also implies this limit in greater
generality.  Our linear-form argument is included because it is shorter
in the fixed-seed orbit and interfaces directly with the radical
identity.  Related dyadic Lucas factorizations and square-free kernels
were studied by Ribenboim~\cite{Ribenboim2002}.

For the illustrative seed \((1,8,9)\), the nonarchimedean terms belong
to the known sequence
\[
 3^mT_m(1/3)=\frac12V_m(2,9),
\]
whose terms and Chebyshev--Lucas description are recorded as OEIS
A025172~\cite{OEIS-A025172}.  Pairwise coprimality at dyadic indices,
rank congruences, and the individual Lucas--Wieferich criterion are
standard consequences of Lucas-sequence
theory~\cite{Ribenboim2001,Sun2013}.  The orbit identity, its telescoped
effective estimate, and the resulting aggregate quality equivalence are
the contribution of this note.  A bounded search of the mathematical
literature and citation trails through July 26, 2026 found no earlier
source containing these whole-orbit statements; this is not a claim that
the search was exhaustive.

The conclusions do not prove or disprove the \(abc\) conjecture.

\section{A general quadratic orbit}

Let
\[
 0<a_0<b_0,\qquad a_0+b_0=c_0,\qquad \gcd(a_0,b_0)=1,
                                                               \tag{2.1}
\]
and suppose that \(a_0,b_0\) have opposite parity.  This ordering loses
no generality because the transfer is symmetric in \(a_0,b_0\).  Define
\[
 (a_{n+1},b_{n+1},c_{n+1})
   =\bigl(4a_nb_n,(a_n-b_n)^2,c_n^2\bigr),\qquad
 d_n=a_n-b_n.                                             \tag{2.2}
\]
Identity (1.1) gives \(a_n+b_n=c_n\), and
\[
 c_n=c_0^{2^n},\qquad d_{n+1}=c_n^2-2d_n^2.               \tag{2.3}
\]

\begin{lemma}[Prime supports]\label{lem:support}
Every triple in (2.2) is primitive.  The integers \(d_j\) are odd,
pairwise coprime, and coprime to \(a_0b_0c_0\).  Consequently, if
\(R_n=\rad(a_nb_nc_n)\), then
\[
 R_n=R_0\prod_{j=0}^{n-1}\rad(|d_j|).                     \tag{2.4}
\]
\end{lemma}

\begin{proof}
Opposite parity is preserved by (2.2), so \(c_n,d_n\) are odd.  At the
seed,
\(\gcd(c_0,d_0)=\gcd(a_0+b_0,a_0-b_0)=1\), since this gcd divides
\(2\gcd(a_0,b_0)\) and both arguments are odd.  If
\(\gcd(c_n,d_n)=1\), then
\[
 \gcd(c_{n+1},d_{n+1})
 =\gcd(c_n^2,c_n^2-2d_n^2)
 =\gcd(c_n^2,2d_n^2)=1,
\]
because \(c_n\) is odd.  Thus \(\gcd(c_n,d_n)=1\) for every \(n\).
Any common divisor of \(a_n,b_n\) divides both \(c_n\) and \(d_n\), so
every triple is primitive.

We first exclude the primes of the seed.  Let \(p\) be an odd prime.
If \(p\mid c_0\), then \(p\mid c_k\) for every \(k\), whereas
\(p\nmid d_0\); (2.3) shows inductively that \(p\nmid d_k\).  If
\(p\mid a_0\), then \(d_0\equiv-c_0\pmod p\), and the congruence
\(d_k\equiv-c_k\pmod p\) propagates under (2.3).  If \(p\mid b_0\),
then \(d_0\equiv c_0\pmod p\), after which
\(d_1\equiv-c_1\pmod p\) and the same invariant propagates.  Thus no
prime of \(a_0b_0c_0\) divides a \(d_j\).

If \(p\mid d_m\), then \(p\nmid c_m\), and (2.3) gives
\[
 d_{m+1}\equiv c_{m+1}\pmod p,\qquad
 d_{m+2}\equiv-c_{m+2}\pmod p.
\]
The latter congruence persists, so \(p\) divides no later \(d_k\).
This proves pairwise coprimality.

For \(n\geq1\), \(b_n=d_{n-1}^2\), while induction in (2.2) shows that
the prime support of \(a_n\) is the support of \(a_0b_0\), together
with the supports of \(d_0,\ldots,d_{n-2}\).  All these supports and
the support of \(c_0\) are disjoint.  Equation (2.4) follows.
\end{proof}

Set
\[
 x_0=\frac{b_0-a_0}{c_0}=\cos\theta,\qquad
 0<\theta<\frac{\pi}{2}.
\]
Normalizing (2.3) gives
\[
 -\frac{d_n}{c_n}=T_{2^n}(x_0)=\cos(2^n\theta).            \tag{2.5}
\]
The complex number
\[
 z=e^{i\theta}
   =\frac{b_0-a_0+2\sqrt{-a_0b_0}}{c_0}                  \tag{2.6}
\]
is not a root of unity.  Indeed, if it were, then
\(z+z^{-1}=2(b_0-a_0)/c_0\) would be a rational algebraic integer.
Since \(\gcd(c_0,b_0-a_0)=1\) and \(c_0\) is odd and greater than \(1\),
this is impossible.  In particular, none of the sine or cosine values
used below vanishes.

Define
\[
 t_j=\frac{|d_j|}{c_j}=|\cos(2^j\theta)|,\qquad
 W_n=\prod_{j=0}^{n-1}\cosoc(|d_j|),\qquad
 \cosoc(m)=\frac{m}{\rad(m)}.                             \tag{2.7}
\]

\begin{theorem}[Exact orbit identity]\label{thm:exact}
For every \(n\geq0\),
\[
 \boxed{\quad
 \frac{R_n}{c_n}
 =\frac{R_0}{c_0}\,
   \frac{|\sin(2^n\theta)|}
        {2^n\sin\theta\,W_n}.
 \quad}                                                    \tag{2.8}
\]
Consequently,
\[
 \frac{c_n}{R_n}
 \geq\frac{2^{n+1}\sqrt{a_0b_0}}{R_0}
 =\frac{2\sqrt{a_0b_0}}{R_0\log c_0}\log c_n.             \tag{2.9}
\]
Thus every orbit satisfying (2.1) eventually consists of \(abc\)-hits.
If the seed is already a hit, then every member of the orbit is a hit.
\end{theorem}

\begin{proof}
Since
\[
 \prod_{j<n}c_j=c_0^{2^n-1}=\frac{c_n}{c_0},
\]
equations (2.4) and (2.7) first give
\[
 \frac{R_n}{c_n}
 =\frac{R_0}{c_0}\frac{\prod_{j<n}t_j}{W_n}.              \tag{2.10}
\]
Iterating \(\sin(2x)=2\sin x\cos x\) gives the exact telescope
\[
 \prod_{j=0}^{n-1}|\cos(2^j\theta)|
 =\frac{|\sin(2^n\theta)|}{2^n\sin\theta},                \tag{2.11}
\]
which proves (2.8).  Now use \(W_n\geq1\),
\(|\sin(2^n\theta)|\leq1\), and
\(\sin\theta=2\sqrt{a_0b_0}/c_0\) to obtain (2.9).
If \(R_0<c_0\), (2.10) and \(0<t_j<1\) prove the final assertion.
\end{proof}

\begin{remark}
The lower bound (2.9) has order only \(\log c_n\).  Much larger
unconditional excesses \(c/\rad(abc)\) are known in specially designed
families.  In particular, Bright~\cite{Bright2024} proves that
\[
 c>\rad(abc)\exp\left(
  \frac{6.563\sqrt{\log c}}{\log\log c}\right)
\]
for infinitely many primitive triples; see also
\cite{GranvilleTucker2002,AlvarezEtAl2023}.  The point of (2.8) is its
exactness and its validity for every primitive opposite-parity seed.
\end{remark}

\section{Effective archimedean control and quality}

We record a completely effective estimate for the only oscillating factor
in (2.8).  Let
\[
 C_{\mathrm{BW}}
 =18\cdot3!\cdot2^3\cdot64^4\log8,\qquad
 H_0=\max\left\{\frac12\log c_0,\frac{\theta}{2},\frac12\right\},
\qquad
 \kappa=C_{\mathrm{BW}}H_0\frac{\pi}{2}.                  \tag{3.1}
\]

\begin{theorem}[One-shot Baker estimate]\label{thm:baker}
For \(n\geq2\),
\[
 |\sin(2^n\theta)|>\frac{2}{\pi}\,2^{-\kappa(n+1)}.        \tag{3.2}
\]
If
\[
 F_n=\log\frac{c_n}{R_n}-\log W_n,\qquad
 C_0=\log\frac{c_0}{R_0}+\log\sin\theta,
\]
then the lower bound below holds for every \(n\geq0\), and the upper
bound holds for \(n\geq2\):
\[
 C_0+n\log2
 \leq F_n
 <C_0+\log\frac{\pi}{2}+\kappa\log2
      +(1+\kappa)n\log2.                                  \tag{3.3}
\]
In particular,
\[
 \log\frac{c_n}{R_n}=\log W_n+\Theta_{a_0,b_0}(n).        \tag{3.4}
\]
\end{theorem}

\begin{proof}
Put \(N=2^n\), choose an integer \(k\) nearest to \(N\theta/\pi\), and
take the principal logarithms \(\log z=i\theta\) and
\(\log(-1)=i\pi\).  Then
\[
 \Lambda_N=N\log z-k\log(-1)
           =i(N\theta-k\pi),\qquad |\Lambda_N|\leq\pi/2.  \tag{3.5}
\]
The form is nonzero because \(z\) is not a root of unity.  Moreover
\(|k|\leq N\).

We apply the explicit linear-form estimate of Baker and
Wüstholz~\cite[Theorem~1]{BakerWustholz1993}.  The field
\(\QQ(z,-1)=\QQ(\sqrt{-a_0b_0})\) has degree \(D=2\).  The primitive
minimal polynomial
\[
 c_0X^2-2(b_0-a_0)X+c_0
\]
has both roots on the unit circle, so \(h(z)=\frac12\log c_0\).
The modified heights in the theorem may therefore be taken as \(H_0\)
and \(\pi/2\).  For \(n\geq2\), take the strict coefficient bound
\(B=2N\geq8\).
The theorem gives
\[
 \log|\Lambda_N|>-\kappa\log(2N).                         \tag{3.6}
\]
Since
\[
 |\sin(N\theta)|
 =|\sin(N\theta-k\pi)|
 \geq\frac{2}{\pi}|\Lambda_N|,
\]
equation (3.2) follows.  Finally, taking logarithms in (2.8) gives the
exact expression
\[
 F_n=C_0+n\log2-\log|\sin(2^n\theta)|.                    \tag{3.7}
\]
The inequalities \(0<|\sin(2^n\theta)|\leq1\) and (3.2) give (3.3).
\end{proof}

\begin{corollary}[Quality criterion]\label{cor:quality}
Put \(q_n=\log c_n/\log R_n\).  Then
\[
 q_n\longrightarrow1
 \quad\Longleftrightarrow\quad
 \log W_n=o(\log c_n),                                    \tag{3.8}
\]
and
\[
 \limsup_{n\to\infty}(q_n-1)>0
 \quad\Longleftrightarrow\quad
 \limsup_{n\to\infty}\frac{\log W_n}{\log c_n}>0.          \tag{3.9}
\]
For every fixed \(\delta>0\), the exact threshold is
\[
 q_n\geq1+\delta
 \quad\Longleftrightarrow\quad
 \log W_n+F_n\geq
 \frac{\delta}{1+\delta}\log c_n.                         \tag{3.10}
\]
\end{corollary}

\begin{proof}
By (2.9), \(R_n<c_n\) for all sufficiently large \(n\).  Write
\(\Delta_n=\log(c_n/R_n)=\log W_n+F_n\).  Since
\(\log c_n=2^n\log c_0\), Theorem~\ref{thm:baker} gives
\(F_n=o(\log c_n)\).  Now
\[
 q_n-1=\frac{\Delta_n}{\log R_n},
\]
which proves (3.8)--(3.9).  Rearranging this identity gives (3.10).
\end{proof}

\begin{proposition}[General Lucas form]\label{prop:general-lucas}
For the Lucas \(V\)-sequence with parameters
\[
 P=2(b_0-a_0),\qquad Q=c_0^2,
\]
which satisfy \(\gcd(P,Q)=1\), one has
\[
 d_j=-\frac12V_{2^j}(P,Q).                                \tag{3.11}
\]
\end{proposition}

\begin{proof}
Because \(c_0\) is odd and \(\gcd(c_0,b_0-a_0)=1\), the parameters are
coprime.
The two roots of \(X^2-PX+Q\) are
\(\gamma=b_0-a_0+2\sqrt{-a_0b_0}\) and its conjugate.  Thus
\[
 \frac12V_m(P,Q)=c_0^mT_m\left(\frac{b_0-a_0}{c_0}\right).
\]
Take \(m=2^j\) and compare with (2.5).
\end{proof}

\begin{corollary}[Dyadic negative-discriminant triples]
\label{cor:negative-discriminant}
Let \(U_m=U_m(P,Q)\) be the companion sequence to the \(V\)-sequence
above, and put
\[
 D=P^2-4Q=-16a_0b_0.
\]
For \(m=2^n\),
\[
 (a_{n+1},b_{n+1},c_{n+1})
 =\left(\frac{|D|U_m^2}{4},\frac{V_m^2}{4},Q^m\right).
                                                               \tag{3.12}
\]
Moreover,
\[
 \lim_{n\to\infty}
 \frac{\log|U_{2^n}^2|}{\log Q^{2^n}}=1.                 \tag{3.13}
\]
\end{corollary}

\begin{proof}
The roots in the proof of Proposition~\ref{prop:general-lucas} are
\(\gamma=c_0e^{i\theta}\) and \(\bar\gamma=c_0e^{-i\theta}\).  Hence
\[
 U_m=\frac{c_0^m\sin(m\theta)}{2\sqrt{a_0b_0}}.
\]
For \(m=2^n\), equations (2.3) and (2.5) give
\[
 a_nb_n=\frac{c_n^2-d_n^2}{4}=a_0b_0U_m^2.
\]
Together with \(V_m/2=-d_n\), this proves (3.12).  Also,
\[
 \log|U_m^2|
 =\log Q^m+2\log|\sin(m\theta)|-\log(4a_0b_0).
\]
Theorem~\ref{thm:baker} makes the last two terms \(O(n)\), while
\(\log Q^{2^n}=2^{n+1}\log c_0\), proving (3.13).
\end{proof}

\begin{remark}
For \(n\geq1\), both \(U_{2^n}\) and \(V_{2^n}\) are even, so (3.12) is
the normalized \(1/4\)-case of the Lucas \(abc\)-triples in
\cite[\S2.2]{Bolvardizadeh2023}.  Bolvardizadeh later singled out the
limit (3.13) as an obstacle for negative discriminants.  Hajdu and
Tijdeman's subsequent all-index growth theorem
\cite[Theorem~2.1]{HajduTijdeman2025} implies that limit for every
nondegenerate non-real Lucas sequence.  The proof above is retained only
as a short dyadic derivation which exposes the same sine factor appearing
in (2.8); no all-index novelty is claimed.  Neither estimate by itself
proves an unconditional quality limit.
\end{remark}

\section{The orbit of (1,8,9)}

We now specialize to
\[
 (a_0,b_0,c_0)=(1,8,9),\qquad
 \theta=\arccos(7/9).
\]
The first triples are
\[
 (1,8,9),\quad(32,49,81),\quad(6272,289,6561).
\]
Since \(R_0=6\) and \(\sin\theta=4\sqrt2/9\), (2.8) becomes
\[
 \boxed{\quad
 \frac{R_n}{c_n}
 =\frac{3\,|\sin(2^n\theta)|}
        {2^{n+1}\sqrt2\,W_n}.
 \quad}                                                    \tag{4.1}
\]

There is a smaller-parameter Lucas description in this case.  Let
\[
 A_m=3^mT_m(1/3)=\frac12V_m(2,9).
\]
The composition identity for Chebyshev polynomials and
\(T_2(1/3)=-7/9\) give
\[
 |d_j|=|A_{2^{j+1}}|
      =\frac12\bigl|V_{2^{j+1}}(2,9)\bigr|.                \tag{4.2}
\]
As noted in the introduction, this sequence identity is prior art
\cite{OEIS-A025172}.

\subsection{Exact Lucas--Wieferich contributions}

Let \(U_m=U_m(2,9)\), and put
\[
 \alpha=1+2\sqrt{-2},\qquad
 \beta=1-2\sqrt{-2},\qquad u=\alpha/\beta.
\]
Then
\[
 U_m=\frac{\alpha^m-\beta^m}{\alpha-\beta},\qquad
 V_m=\alpha^m+\beta^m.
\]

\begin{proposition}[Rank and valuation]\label{prop:valuation}
Suppose that \(p\mid d_j\), put \(m=2^{j+1}\), and let
\(\chi_p=\left(\frac{-2}{p}\right)\).  Then \(p\nmid6\).  For a prime
\(\mathfrak p\mid p\), write \(\widetilde u\) for the image of \(u\) in
\((\OK/\mathfrak p)^\times\).  Then
\[
 \ord_{(\OK/\mathfrak p)^\times}(\widetilde u)=2m=2^{j+2}
 \quad\text{for every prime }\mathfrak p\mid p
 \text{ in }\QQ(\sqrt{-2}).                               \tag{4.3}
\]
Consequently,
\[
 2^{j+2}\mid p-\chi_p.                                    \tag{4.4}
\]
Moreover,
\[
 v_p(d_j)=v_p\bigl(U_{p-\chi_p}(2,9)\bigr).                \tag{4.5}
\]
In particular, \(p^2\mid d_j\) if and only if \(p\) is a
Lucas--Wieferich prime for \(U(2,9)\).
\end{proposition}

\begin{proof}
Lemma~\ref{lem:support} gives \(p\nmid18\), and \(d_j\) is odd, so
\(p\nmid6\).  The rank statements are standard; we include the short
argument.
Modulo \(\mathfrak p\), (4.2) gives \(u^m=-1\).  Since \(m\) is a power
of \(2\), the order is exactly \(2m\).  If \(p\) splits, this order
divides \(p-1\); if \(p\) is inert, Frobenius gives
\(\widetilde u^{\,p}=\widetilde{\,\overline u\,}
 =\widetilde u^{-1}\), so the order divides
\(p+1\).  This proves (4.4), equivalently the rank criterion in
Ribenboim~\cite[\S2.13]{Ribenboim2001}.

Sun's Lucas law of repetition
\cite[Theorem~3(ii)]{Sun2013} states, under the hypotheses satisfied by
\((P,Q)=(2,9)\), that an odd prime \(p\mid V_m(P,Q)\) satisfies
\[
 v_p(V_m)=v_p(m)+v_p\left(U_{p-(D/p)}(P,Q)\right),
 \qquad D=P^2-4Q.
\]
Here \(D=-32\), \((D/p)=\chi_p\), and \(v_p(m)=0\).  Division by \(2\)
does not affect the \(p\)-adic valuation, proving (4.5).
\end{proof}

\begin{remark}
Sun's theorem and his analogous iterated companions
\(S_{k+1}=S_k^2-2\) already provide the per-prime valuation analysis
used in Proposition~\ref{prop:valuation}.  Our new statement is the
coupling of all such valuations across the \(abc\)-orbit in (4.6), not
the individual Lucas--Wieferich criterion.
\end{remark}

Combining Proposition~\ref{prop:valuation} with pairwise coprimality
gives the exact expansion
\[
 \boxed{\quad
 W_n=\prod_{j=0}^{n-1}\ \prod_{p\mid d_j}
 p^{\,v_p(U_{p-\chi_p}(2,9))-1}.
 \quad}                                                    \tag{4.6}
\]
Thus (3.8)--(3.9) are precisely statements about the aggregate weight of
Lucas--Wieferich primes encountered at dyadic ranks.  This is more
specific than the classical one-directional implications from the
\(abc\) conjecture to the existence of non-Wieferich primes
\cite{Silverman1988,RibenboimWalsh1999,Yabuta2007}.

\subsection{What unconditional radical bounds give}

There are unconditional results directly relevant to (4.2).
Stewart~\cite[Theorem~1]{Stewart1983} proves an effective radical lower
bound for Lucas and Lehmer sequences.  Applied to the companion sequence
\(V_m(2,9)\) at \(m=2^k\), it yields, for all sufficiently large \(k\),
\[
 \log\rad\bigl(|V_{2^k}(2,9)|\bigr)
 \gg \frac{k^2}{\log k}.                                  \tag{4.7}
\]
Indeed, in Stewart's notation \(q(m)\) is the number of square-free
divisors of \(m\), \(d(m)\) is the number of positive divisors, and
\(|2|_2=1/2\).  Thus for \(m=2^k\), one has \(q(m)=2\) and
\(d(m|m|_2)=1\), so the \(v_m\)-sequence clause of his theorem has
precisely the order of growth in (4.7).
For the roots \(\alpha,\beta\) above,
\((\alpha+\beta)^2=4\), \(\alpha\beta=9\), and
\(\gcd(4,9)=1\); moreover \(\alpha/\beta\) is not a root of unity
because \(u+u^{-1}=-14/9\).  These verify the hypotheses of Stewart's
companion-sequence estimate.
Since removing the factor \(2\) changes the logarithm of the radical by
at most \(\log2\), summation in (2.4) gives
\[
 \log R_n\gg\frac{n^3}{\log n},\qquad
 q_n\ll\frac{2^n\log n}{n^3}.                             \tag{4.8}
\]
This is far below the exponential-in-\(n\) radical scale required by
(3.8), but it explains precisely what presently available general
theorems contribute.  Primitive-divisor results alone address the
appearance of new primes, whereas here every prime is already confined
to one dyadic term by Lemma~\ref{lem:support}; their multiplicities are
the unresolved issue; compare Bilu, Hanrot, and
Voutier~\cite{BiluHanrotVoutier2001}.

\section{Computation}

The accompanying program \texttt{check\_square\_lifts.py} works entirely
modulo \(p^2\).  For each prime \(p\), it iterates
\[
 (c,d)\longmapsto(c^2,c^2-2d^2)\pmod{p^2},
 \qquad(c,d)=(9,-7),
\]
and reports \(p,j\) exactly when \(d_j\equiv0\pmod{p^2}\).  If
\(p\mid d_j\) but \(p^2\nmid d_j\), the search for that prime stops:
Lemma~\ref{lem:support} proves that \(p\) cannot occur later.

\begin{proposition}[Finite verification]\label{prop:computation}
For every prime \(p\leq10^7\), \(p\neq2,3\), and every
\(0\leq j\leq50\), the modular recurrence gives \(p^2\nmid d_j\).
The search tests \(664{,}577\) primes and reports zero square lifts.
\end{proposition}

This finite verification is not an input to any theorem.  The script
contains a deterministic exact-versus-modular self-check for small
indices.  The command used for Proposition~\ref{prop:computation} is
\begin{center}
\verb|python3 check_square_lifts.py --prime-limit 10000000 --max-j 50|.
\end{center}
The first terms provide short independent checks of the recurrence,
factorization, and congruence (4.4):
\[
\begin{array}{c@{\quad}r@{\quad}l@{\quad}l}
\toprule
j&d_j&|d_j|\text{ factored}&p\bmod 2^{j+2}\\
\midrule
0&-7&7&-1\\
1&-17&17&+1\\
2&5983&31\cdot193&-1,+1\\
3&-28545857&2753\cdot10369&+1,+1\\
4&223288285122943&127\cdot19841\cdot88613249&-1,+1,+1\\
\bottomrule
\end{array}
\]

\section*{Code and data availability}

The source package accompanying this article contains
\texttt{check\_square\_lifts.py}.  It uses only the Python~3 standard
library and generates the reported counts directly; no external data
set is used.  The command and all search bounds are stated above.

\section{Conclusion}

The exact identity (2.8) isolates the two possible sources of radical
loss in every primitive opposite-parity quadratic-transfer orbit.  The
archimedean source is completely controlled: it contributes only
\(\Theta(n)\) to \(\log(c_n/R_n)\).  Any fixed asymptotic quality gap
must therefore come from positive-power accumulation in \(W_n\).

For the seed \((1,8,9)\), formula (4.6) identifies that accumulation
exactly with weighted Lucas--Wieferich valuations at dyadic ranks.  The
finite search and the known radical bounds are consistent with
sub-power growth, but neither decides it.  Proving
\[
 \sum_{j<n}\log\frac{|V_{2^{j+1}}(2,9)|/2}
                         {\rad(|V_{2^{j+1}}(2,9)|/2)}
 =o(2^n)
\]
would prove \(q_n\to1\) for this orbit; positive linear growth on the
\(2^n\) scale along a subsequence would instead give a fixed quality
gap.  Both remain open.

\section*{AI-use statement}

Anthropic Claude Fable~5 and OpenAI Codex (GPT-5-based service model),
using the service versions available on July 26, 2026, were used for
idea exploration, candidate proofs, literature classification and
retrieval, coding assistance, drafting, and reciprocal adversarial
review.  They were used to accelerate exploration and to expose errors
through independent cross-checking.  Mathematical claims were
line-checked in separate review passes, bibliographic claims were
checked against cited sources, and the computation was independently
reproduced.  The AI systems are not authors.  The named author directed
the project and accepts responsibility for the originality, validity,
references, code, and final content.

\begin{thebibliography}{99}

\bibitem{AlvarezEtAl2023}
E.~Alvarez-Salazar, A.~J.~Barrios, C.~Henaku, and S.~Soller,
\newblock On \(abc\) triples of the form \((1,c-1,c)\),
\newblock \emph{Integers} \textbf{23} (2023), Paper A64, 22 pp.
\newblock \url{https://doi.org/10.5281/zenodo.8283159}.

\bibitem{BakerWustholz1993}
A.~Baker and G.~W{\"u}stholz,
\newblock Logarithmic forms and group varieties,
\newblock \emph{J. Reine Angew. Math.} \textbf{442} (1993), 19--62.
\newblock \url{https://doi.org/10.1515/crll.1993.442.19}.

\bibitem{BiluHanrotVoutier2001}
Y.~Bilu, G.~Hanrot, and P.~M.~Voutier,
\newblock Existence of primitive divisors of Lucas and Lehmer numbers,
\newblock \emph{J. Reine Angew. Math.} \textbf{539} (2001), 75--122.
\newblock \url{https://doi.org/10.1515/crll.2001.080}.

\bibitem{Bolvardizadeh2023}
S.~Bolvardizadeh,
\newblock \emph{On the quality of the \(ABC\)-solutions},
\newblock M.Sc. thesis, University of Lethbridge, 2023.
\newblock \url{https://hdl.handle.net/10133/6591}.

\bibitem{Bright2024}
C.~Bright,
\newblock A new lower bound in the \(abc\) conjecture,
\newblock \emph{Canad. Math. Bull.} \textbf{67} (2024), 369--378.
\newblock \url{https://doi.org/10.4153/S0008439523000784}.

\bibitem{GranvilleTucker2002}
A.~Granville and T.~J.~Tucker,
\newblock It's as easy as \(abc\),
\newblock \emph{Notices Amer. Math. Soc.} \textbf{49} (2002),
1224--1231.
\newblock \url{https://www.ams.org/notices/200210/fea-granville.pdf}.

\bibitem{Guninski2020}
G.~Guninski (posting as \texttt{joro}),
\newblock Answer to ``How balanced can \(abc\) triples be?'',
\newblock MathOverflow, April 2, 2020.
\newblock \url{https://mathoverflow.net/questions/356295/},
accessed July 26, 2026.

\bibitem{HajduTijdeman2025}
L.~Hajdu and R.~Tijdeman,
\newblock Integers represented by Lucas sequences,
\newblock \emph{Ramanujan J.} \textbf{66} (2025), Article 74, 27 pp.
\newblock \url{https://doi.org/10.1007/s11139-025-01041-6}.

\bibitem{Holsztynski2017}
W{\l}odzimierz Holszty\'nski,
\newblock Answer to ``\(abc\) streams (sequences of creek stones)'',
\newblock MathOverflow, March 15, 2017.
\newblock \url{https://mathoverflow.net/questions/263463/},
accessed July 26, 2026.

\bibitem{MartinMiao2016}
G.~Martin and W.~Miao,
\newblock \(abc\) triples,
\newblock \emph{Funct. Approx. Comment. Math.} \textbf{55} (2016),
145--176.
\newblock \url{https://doi.org/10.7169/facm/2016.55.2.2}.

\bibitem{OEIS-A025172}
N.~J.~A.~Sloane and The OEIS Foundation,
\newblock The On-Line Encyclopedia of Integer Sequences, A025172,
\newblock \url{https://oeis.org/A025172}, accessed July 26, 2026.

\bibitem{Oesterle1988}
J.~Oesterl\'e,
\newblock Nouvelles approches du ``th\'eor\`eme'' de Fermat,
\newblock S\'eminaire Bourbaki, Exp.\ 694,
\emph{Ast\'erisque} \textbf{161--162} (1988), 165--186.
\newblock \url{https://www.numdam.org/item/SB_1987-1988__30__165_0/}.

\bibitem{OhanaSpicerStein2013}
R.~A.~Ohana, S.~Spicer, and W.~Stein,
\newblock \emph{The \(ABC\) data},
\newblock unpublished manuscript, October 2013, 13 pp.
\newblock
\href{https://cocalc.com/share/download/c1f4c5685b89bae0dfa24156574398b8c8172a3a/briefing/brief.pdf}
{Online manuscript}, accessed July 26, 2026.

\bibitem{Ribenboim2001}
P.~Ribenboim,
\newblock On square factors of terms of binary recurring sequences and
the \(ABC\) conjecture,
\newblock \emph{Publ. Math. Debrecen} \textbf{59} (2001), 459--469.
\newblock \url{https://doi.org/10.5486/PMD.2001.2559}.
\newblock Full text:
\url{https://publi.math.unideb.hu/paper/752/download/10_5486_PMD_2001_2559.pdf}.

\bibitem{Ribenboim2002}
P.~Ribenboim,
\newblock The square-free kernel of \(x^{2^n}-a^{2^n}\),
\newblock \emph{Acta Arith.} \textbf{101} (2002), 189--197.
\newblock \url{https://doi.org/10.4064/aa101-2-9}.

\bibitem{RibenboimWalsh1999}
P.~Ribenboim and P.~G.~Walsh,
\newblock The \(ABC\) conjecture and the powerful part of terms in binary
recurring sequences,
\newblock \emph{J. Number Theory} \textbf{74} (1999), 134--147.
\newblock \url{https://doi.org/10.1006/jnth.1998.2315}.

\bibitem{Silverman1988}
J.~H.~Silverman,
\newblock Wieferich's criterion and the \(abc\)-conjecture,
\newblock \emph{J. Number Theory} \textbf{30} (1988), 226--237.
\newblock \url{https://doi.org/10.1016/0022-314X(88)90019-4}.

\bibitem{Stewart1983}
C.~L.~Stewart,
\newblock On divisors of Fermat, Fibonacci, Lucas and Lehmer numbers, III,
\newblock \emph{J. London Math. Soc.} (2) \textbf{28} (1983), 211--217.
\newblock \url{https://doi.org/10.1112/jlms/s2-28.2.211}.
\newblock Author-hosted scan:
\url{https://uwaterloo.ca/pure-mathematics/sites/default/files/uploads/documents/j-london-math-soc-1983.pdf}.

\bibitem{Sun2013}
Z.-H.~Sun,
\newblock Congruences concerning Lucas' law of repetition,
\newblock arXiv:1312.3511 [math.NT], 2013.
\newblock \url{https://arxiv.org/abs/1312.3511}.

\bibitem{vanderHorst2010}
J.~P.~van der Horst,
\newblock \emph{Finding \(ABC\)-triples using elliptic curves},
\newblock M.Sc. thesis, Universiteit Leiden, 2010.
\newblock
\url{https://math.leidenuniv.nl/scripties/vanderHorstMaster.pdf}.

\bibitem{Yabuta2007}
M.~Yabuta,
\newblock The \(ABC\)-conjecture and the powerful numbers in Lucas
sequences,
\newblock \emph{Fibonacci Quart.} \textbf{45} (2007), 362--365.
\newblock \url{https://doi.org/10.1080/00150517.2007.12428206}.

\end{thebibliography}

\end{document}
